\subsection{Multimeter}
\markright{Multimeter}

\subsubsection*{Definition}
A multimeter---also called a volt-ohm meter (VOM)---is a multi-function electronic
measuring instrument that combines the capabilities of a voltmeter, an ammeter, and an
ohmmeter in a single portable unit \cite{wiki:multimeter}. Modern digital multimeters
(DMMs) convert measured quantities to digital form and display them on an LCD. Older
analog multimeters deflect a pointer over a calibrated scale using a moving-coil
D'Arsonval galvanometer mechanism.

\labimage{lab-01/assets/multimeter.jpg}%
  {A digital multimeter (DMM): the central rotary dial selects the
   measurement function and range; probes connect via colour-coded banana sockets.}%
  {multimeter}

\subsubsection*{Specifications}
A typical handheld DMM measures DC voltage (200\,mV to 600\,V), AC voltage (up to
600\,V), DC current (200\,$\mu$A to 10\,A), and resistance (200\,$\Omega$ to
200\,M$\Omega$) \cite{electronicsclub:multimeter}. Input impedance is standardized at
10\,M$\Omega$ to minimize circuit loading. DC voltage accuracy for a quality DMM is
typically $\pm(0.5\%\,\mathrm{rdg} + 1\,\mathrm{digit})$; laboratory bench instruments
achieve better than $\pm0.01\%$ of reading. Many DMMs also measure capacitance,
frequency, diode forward voltage, and transistor $h_{FE}$.

\subsubsection*{Purpose of Use}
The multimeter verifies operating voltages, checks component values, tests continuity
(detecting broken or short-circuit connections), measures the forward voltage of diodes,
and confirms the polarity of electrolytic capacitors. As a voltmeter it is connected
in parallel; as an ammeter it must be inserted in series via the dedicated current
terminal.

\subsubsection*{Applications}
Multimeters are indispensable in electronics laboratories for circuit debugging and
verification, in field service for diagnosing household wiring, automotive electrical
faults, and industrial control panels \cite{wiki:multimeter}. Clamp-meter variants
extend the same measurement philosophy to high AC currents without breaking the circuit.
